% Generated draft: requires mathematical and source review.
\documentclass[11pt]{article}
\usepackage[margin=25mm]{geometry}
\usepackage{fontspec}
\setmainfont{Latin Modern Roman}
\usepackage{amsmath,amssymb,mathtools,unicode-math}
\setmathfont{Latin Modern Math}
\usepackage{enumitem,needspace,xurl,hyperref}
\setlength{\parindent}{0pt}
\setlength{\parskip}{5pt}
\setlength{\emergencystretch}{3em}
\newcommand{\N}{\mathbb{N}}
\newcommand{\Z}{\mathbb{Z}}
\newcommand{\Q}{\mathbb{Q}}
\newcommand{\R}{\mathbb{R}}
\newcommand{\C}{\mathbb{C}}
\newcommand{\F}{\mathbb{F}}
\newcommand{\A}{\mathbb{A}}
\newcommand{\PP}{\mathbb{P}}
\newcommand{\E}{\mathbb{E}}
\newcommand{\eps}{\varepsilon}
\newcommand{\dd}{\,\mathrm{d}}
\begin{document}
\section*{UnsolvedMath ID: 800807}
\textbf{Source problem code:} catalogue entry
\subsection{Definitions and mathematical statement}
The source asks for a physical or fluid model realizing an \(L^2\) left-invariant metric on the volume-preserving diffeomorphism group \(\mathscr{D}_\mu(M)\). The geometry of \(M\), boundary conditions, and regularity class are not specified. The following describes the metric condition conditionally, without selecting those missing hypotheses.

Given a smooth Riemannian manifold \((M,g)\) and a smooth positive volume density \(\mu\), write
\[
\mathscr{D}_\mu(M)=\{\varphi:M\to M:\ \varphi\text{ is a diffeomorphism and }\varphi^*\mu=\mu\},
\]
with whatever regularity and boundary restrictions are intended. At the identity, the \(L^2\) inner product on admissible divergence-free vector fields is
\[
\langle u,v\rangle_{L^2}=\int_M g_x(u(x),v(x))\,\dd\mu(x),
\]
provided the integral is finite. Its left-invariant extension is characterized by
\[
G_\varphi(U,V)=\int_M g_x\!\left((D\varphi_x)^{-1}U(x),(D\varphi_x)^{-1}V(x)\right)\,\dd\mu(x),
\]
where \(U,V\in T_\varphi\mathscr{D}_\mu(M)\) are vector fields along \(\varphi\). Equivalently, with \(L_\eta(\varphi)=\eta\circ\varphi\), left invariance means
\[
G_{\eta\circ\varphi}\bigl(T_\varphi L_\eta U,T_\varphi L_\eta V\bigr)=G_\varphi(U,V)
\]
for every admissible \(\eta,\varphi,U,V\).

\textbf{Problem.} Describe a physical or fluid model whose configuration space is \(\mathscr{D}_\mu(M)\) and whose kinetic-energy metric has this \(L^2\) left-invariant form. Identify physically interesting systems exhibiting this invariance. The supplied text does not specify a criterion for an acceptable model, nor whether the request concerns one choice of \(M\) or a general construction.

The supplied context also requests investigation of an unspecified preceding system for generic \(S\) and \(\rho(\omega)\), and the definition and study of a model with configuration space \(SL(2)\), described as more ``liquid-like.'' It states that in ``all other cases'' the preceding metric is neither left- nor right-invariant. These contextual requests cannot be formulated precisely from this record: the preceding system, the exceptional cases, \(S\), \(\rho\), \(\omega\), the meaning of genericity, and the scalar field of \(SL(2)\) are not supplied.
\subsection{Short English statement}
Find a physically meaningful model on the group of volume-preserving diffeomorphisms whose kinetic energy comes from an \(L^2\) inner product transported by left translations. The source also suggests related investigations involving unspecified parameters and \(SL(2)\). The geometric setting, boundary conditions, and criteria for a physical model remain unresolved.
\subsection{Sources}
\begin{itemize}
\item Local UnsolvedMath registry, permanent ID \(800807\), exact source problem code \texttt{catalogue entry}, title \emph{Geometric Hydrodynamics in Open Problems --- Problem 6}. Supplied upstream revision: \texttt{d371ee917ac9647fbb06cd6fc529deec6f3dda82}.
\item Khesin, Misiolek, and Modin, \emph{Geometric Hydrodynamics in Open Problems} (2022), Problem 6. Supplied source URL: \url{https://arxiv.org/abs/2205.01143}. The registry attributes the problem to these authors and gives proposed year 2022.
\item Supplied provenance reports access on 2026-10-16 and extraction by \texttt{arxiv\_latex}. These are source-record provenance claims; no independent source visit or verification is asserted here. The registry labels the problem \emph{open}; this label is supplied metadata, not an independently verified literature assessment.
\end{itemize}
\end{document}
